% ══════════════════════════════════════════════════════════════════════
%  Unemployment Rate Forecasting — Republic of Kazakhstan
%  Diploma Pre-Defense Report
% ══════════════════════════════════════════════════════════════════════
\documentclass[12pt,a4paper]{article}

% ── Packages ─────────────────────────────────────────────────────────
\usepackage[utf8]{inputenc}
\usepackage[T2A]{fontenc}
\usepackage[english]{babel}
\usepackage{geometry}
\geometry{margin=2.5cm}

\usepackage{amsmath, amssymb, amsfonts}
\usepackage{booktabs}
\usepackage{longtable}
\usepackage{graphicx}
\usepackage{float}
\usepackage{caption}
\usepackage{subcaption}
\usepackage{hyperref}
\usepackage{xcolor}
\usepackage{pgfplots}
\pgfplotsset{compat=1.18}
\usepackage{pgfplotstable}
\usepackage{tikz}
\usepackage{enumitem}
\usepackage{multirow}
\usepackage{array}
\usepackage[numbers]{natbib}

\hypersetup{
    colorlinks=true,
    linkcolor=blue!60!black,
    citecolor=green!50!black,
    urlcolor=blue!70!black,
}

% ── Custom commands ──────────────────────────────────────────────────
\newcommand{\E}{\mathbb{E}}
\newcommand{\Var}{\operatorname{Var}}
\newcommand{\bx}{\mathbf{x}}
\newcommand{\by}{\mathbf{y}}
\newcommand{\bw}{\mathbf{w}}
\DeclareMathOperator*{\argmin}{arg\,min}

% ══════════════════════════════════════════════════════════════════════
\title{
    \textbf{Forecasting Monthly Unemployment Rate\\
    in the Republic of Kazakhstan}\\[0.5em]
    \large A Comparative Study of Econometric and Machine Learning Approaches\\[0.3em]
    \normalsize Diploma Pre-Defense Report
}
\author{Student Name \\ Department of Data Science}
\date{\today}

\begin{document}
\maketitle
\tableofcontents
\newpage

% ══════════════════════════════════════════════════════════════════════
\section{Introduction}
% ══════════════════════════════════════════════════════════════════════

This study develops a comparative forecasting framework for the monthly
unemployment rate of the Republic of Kazakhstan over a 24-month horizon
(January 2024 -- December 2025). We evaluate eight models spanning
classical econometrics, regularised regression, ensemble learning, and
deep recurrent networks, using a unified evaluation protocol based on
time series cross-validation.

The target variable is the registered unemployment rate (percentage of
the economically active population), sourced from the Bureau of National
Statistics of Kazakhstan (\texttt{stat.gov.kz}).

\paragraph{Research objectives.}
\begin{enumerate}[nosep]
    \item Identify macroeconomic drivers of Kazakhstan's unemployment dynamics.
    \item Compare the predictive accuracy of eight forecasting models.
    \item Quantify uncertainty and assess model robustness via $k$-fold
          time series cross-validation.
\end{enumerate}


% ══════════════════════════════════════════════════════════════════════
\section{Data Description}
% ══════════════════════════════════════════════════════════════════════

\subsection{Dataset Overview}

The primary dataset contains \textbf{192 monthly observations}
(January 2010 -- December 2025) with \textbf{zero missing values} after
compilation. Features include macroeconomic indicators from the National
Bank of Kazakhstan, FRED, and the Bureau of National Statistics.

\begin{table}[H]
\centering
\caption{Feature summary after preprocessing.}
\label{tab:features}
\begin{tabular}{@{}lll@{}}
\toprule
\textbf{Feature} & \textbf{Source} & \textbf{Transformation} \\
\midrule
$y_{t-1}$, $y_{t-12}$      & Target            & Autoregressive lags \\
\texttt{usd\_kzt\_diff}     & National Bank      & 1st difference \\
\texttt{Oil\_crude\_brent\_diff} & FRED          & 1st difference \\
\texttt{Base\_Rate\_diff}   & National Bank      & 1st difference \\
\texttt{gold\_price\_usd\_avg\_diff} & Market    & 1st difference \\
\texttt{rub\_kzt\_diff}     & Derived (USD/RUB)  & 1st difference \\
\texttt{GDP\_growth}        & stat.gov.kz        & Level (\%) \\
\texttt{interest\_rate}     & National Bank      & Level \\
\texttt{total\_cpi\_yoy}    & stat.gov.kz        & Level (\%) \\
\texttt{usd\_kzt\_volatility} & Derived (daily $\sigma$) & Monthly std \\
\bottomrule
\end{tabular}
\end{table}

\subsection{Train--Test Split}

\begin{itemize}[nosep]
    \item \textbf{Training:} January 2011 -- December 2023 ($n \approx 108$ months after differencing/lagging).
    \item \textbf{Test:} January 2024 -- December 2025 ($n = 24$ months).
\end{itemize}

This split respects temporal ordering and avoids data leakage.


% ══════════════════════════════════════════════════════════════════════
\section{Exploratory Data Analysis}
\label{sec:eda}
% ══════════════════════════════════════════════════════════════════════

\subsection{Time Series Characteristics}

\begin{figure}[H]
    \centering
    \includegraphics[width=\textwidth]{figures/02a_unemployment_timeseries.png}
    \caption{Monthly unemployment rate with the train/test boundary (red dashed line).}
    \label{fig:ts}
\end{figure}

The unemployment rate exhibits a strong \textbf{downward trend} from
approximately 6.6\% in 2010 to below 4.8\% by 2025, punctuated by
mild seasonal oscillations.

\subsection{Seasonal Decomposition}

\begin{figure}[H]
    \centering
    \includegraphics[width=\textwidth]{figures/02c_seasonal_decomposition.png}
    \caption{Additive seasonal decomposition (period = 12 months).
             Components: Observed, Trend, Seasonal, Residual.}
    \label{fig:decomp}
\end{figure}

The seasonal component reveals a recurring annual pattern with amplitude
$\approx 0.05$ percentage points --- modest relative to the secular trend.

\subsection{Stationarity Testing}

We apply the \textbf{Augmented Dickey--Fuller} (ADF) test under the null
hypothesis $H_0$: the series has a unit root (non-stationary).

$$
\Delta y_t = \alpha + \beta t + \gamma y_{t-1}
              + \sum_{i=1}^{p} \delta_i \Delta y_{t-i} + \varepsilon_t
$$

\begin{table}[H]
\centering
\caption{ADF test results ($\alpha = 0.05$).
         The test is applied to raw levels; differenced series are stationary
         by construction.}
\label{tab:adf}
\begin{tabular}{@{}lccl@{}}
\toprule
\textbf{Variable} & \textbf{ADF Statistic} & \textbf{$p$-value} & \textbf{Stationary?} \\
\midrule
\texttt{total\_cpi\_yoy}       & $-2.2646$ & $0.1836$ & No  \\
\texttt{gold\_price\_usd\_avg} & $+4.7012$ & $1.0000$ & No  \\
\texttt{usd\_kzt}              & $-0.4190$ & $0.9069$ & No  \\
\texttt{interest\_rate}        & $-1.1099$ & $0.7111$ & No  \\
\texttt{Oil\_crude\_brent}     & $-1.8971$ & $0.3334$ & No  \\
\texttt{GDP\_growth}           & $-2.0566$ & $0.2623$ & No  \\
\texttt{unemployed\_rate}      & $\mathbf{-5.7755}$ & $\mathbf{0.0000}$ & \textbf{Yes} \\
\texttt{Base\_Rate}            & $-0.7437$ & $0.8351$ & No  \\
\texttt{rub\_kzt}              & $+0.7981$ & $0.9916$ & No  \\
\bottomrule
\end{tabular}

\subsection{Correlation Analysis}

\begin{figure}[H]
    \centering
    \includegraphics[width=0.85\textwidth]{figures/02e_correlation_matrix.png}
    \caption{Lower-triangular correlation heatmap of all features.}
    \label{fig:corr}
\end{figure}

\subsection{ACF and PACF}

\begin{figure}[H]
    \centering
    \includegraphics[width=\textwidth]{figures/02f_acf_pacf.png}
    \caption{Autocorrelation (left) and partial autocorrelation (right) functions
             for the unemployment rate. The slow decay of ACF confirms non-stationarity;
             PACF suggests AR(2) structure after differencing.}
    \label{fig:acf}
\end{figure}


% ══════════════════════════════════════════════════════════════════════
\section{Methodology}
\label{sec:method}
% ══════════════════════════════════════════════════════════════════════

\subsection{Preprocessing}

\paragraph{Missing values.}
Time-based interpolation is used to fill any gaps, preserving temporal
continuity. This is preferred over K-NN imputation because K-NN treats
observations as exchangeable, violating the sequential dependence
inherent in time series data.

\paragraph{Feature engineering.}
\begin{itemize}[nosep]
    \item \textbf{Lagged target:} $y_{t-1}$ and $y_{t-12}$ capture
          short-term momentum and annual seasonality.
    \item \textbf{First differencing:} Non-stationary regressors
          (exchange rates, oil, base rate) are differenced:
          $\Delta x_t = x_t - x_{t-1}$.
    \item \textbf{Volatility:} Monthly standard deviation of daily
          USD/KZT rates captures exchange rate uncertainty.
\end{itemize}

\paragraph{Scaling.}
All features are standardised ($\mu = 0$, $\sigma = 1$) using
\texttt{StandardScaler} fitted \emph{only} on training data to prevent
information leakage.

\subsection{Models}

We compare \textbf{eight models} across four families:

\subsubsection{Baselines}

\paragraph{1. Naïve Seasonal Forecast.}
$$\hat{y}_t = y_{t-12}$$
The simplest benchmark: the forecast equals the value from the same
month one year prior.

\paragraph{2. Linear Regression (OLS).}
$$\hat{y}_t = \bx_t^\top \hat{\boldsymbol{\beta}} + \hat{\beta}_0, \qquad
  \hat{\boldsymbol{\beta}} = \argmin_{\boldsymbol{\beta}}
  \sum_{t=1}^{T} (y_t - \bx_t^\top \boldsymbol{\beta})^2$$
with Heteroscedasticity-Consistent (HC1) standard errors.

\subsubsection{Classical Forecasting}

\paragraph{3. SARIMA$(2,1,1)(1,1,1)_{12}$.}
The Seasonal ARIMA model combines autoregressive, differencing, and
moving average components at both regular and seasonal frequencies:
$$\Phi_P(B^s)\,\phi_p(B)\,(1-B)^d(1-B^s)^D y_t
  = \Theta_Q(B^s)\,\theta_q(B)\,\varepsilon_t$$
where $B$ is the backshift operator and $s=12$.

\paragraph{4. SARIMAX.}
Extends SARIMA with exogenous regressors $\bx_t$:
$$\Phi_P(B^s)\,\phi_p(B)\,(1-B)^d(1-B^s)^D y_t
  = \boldsymbol{\beta}^\top \bx_t + \Theta_Q(B^s)\,\theta_q(B)\,\varepsilon_t$$

\subsubsection{Regularised \& Ensemble Methods}

\paragraph{5. Elastic Net.}
Combines $\ell_1$ and $\ell_2$ penalties:
$$\hat{\boldsymbol{\beta}} = \argmin_{\boldsymbol{\beta}}
  \left\{ \frac{1}{2n}\|\by - X\boldsymbol{\beta}\|_2^2
  + \alpha\left[\rho\|\boldsymbol{\beta}\|_1
  + \frac{1-\rho}{2}\|\boldsymbol{\beta}\|_2^2\right] \right\}$$
Hyperparameters $\alpha$ and $\rho$ are selected by time series
cross-validation.

\paragraph{6. Random Forest.}
An ensemble of $B$ regression trees:
$$\hat{y}(\bx) = \frac{1}{B}\sum_{b=1}^{B} T_b(\bx)$$
Hyperparameters (depth, $B$, leaf size) are tuned via
\texttt{RandomizedSearchCV} with \texttt{TimeSeriesSplit}.

\paragraph{7. Gradient Boosting.}
Sequentially fits regression trees to the residuals:
$$\hat{y}^{(m)}(\bx) = \hat{y}^{(m-1)}(\bx) + \eta\,h_m(\bx)$$
where $\eta$ is the learning rate and $h_m$ is the $m$-th base learner
minimising a squared-error loss.

\subsubsection{Neural Networks}

\paragraph{8. LSTM (Long Short-Term Memory).}
A recurrent architecture with gated memory cells:
\begin{align}
    f_t &= \sigma(W_f [\mathbf{h}_{t-1}, \bx_t] + b_f) \label{eq:forget}\\
    i_t &= \sigma(W_i [\mathbf{h}_{t-1}, \bx_t] + b_i) \\
    \tilde{C}_t &= \tanh(W_C [\mathbf{h}_{t-1}, \bx_t] + b_C) \\
    C_t &= f_t \odot C_{t-1} + i_t \odot \tilde{C}_t \\
    o_t &= \sigma(W_o [\mathbf{h}_{t-1}, \bx_t] + b_o) \\
    \mathbf{h}_t &= o_t \odot \tanh(C_t)
\end{align}
Architecture: LSTM(64) $\to$ Dropout(0.2) $\to$ LSTM(32) $\to$
Dropout(0.2) $\to$ Dense(16, ReLU) $\to$ Dense(1).
Lookback window: $L = 12$ months.

\subsection{Evaluation Protocol}

\paragraph{Cross-validation.}
We use \textbf{TimeSeriesSplit} with $k=5$ folds to preserve temporal
ordering.

\paragraph{Metrics.}
\begin{align}
    \text{MAE}   &= \frac{1}{n}\sum_{t=1}^{n} |y_t - \hat{y}_t| \\
    \text{RMSE}  &= \sqrt{\frac{1}{n}\sum_{t=1}^{n} (y_t - \hat{y}_t)^2} \\
    \text{MAPE}  &= \frac{100}{n}\sum_{t=1}^{n}
                    \left|\frac{y_t - \hat{y}_t}{y_t}\right| \\
    \text{sMAPE} &= \frac{200}{n}\sum_{t=1}^{n}
                    \frac{|y_t - \hat{y}_t|}{|y_t| + |\hat{y}_t|} \\
    \text{MASE}  &= \frac{\text{MAE}}{\frac{1}{n-m}\sum_{t=m+1}^{n}
                    |y_t - y_{t-m}|}
                    \quad (m = 12)
\end{align}

\paragraph{Hyperparameter tuning.}
\texttt{RandomizedSearchCV} with 40 iterations is used for Random Forest
and Gradient Boosting. Elastic Net uses \texttt{ElasticNetCV} with a
grid over $\alpha$ and $\rho$.


% ══════════════════════════════════════════════════════════════════════
\section{Results}
\label{sec:results}
% ══════════════════════════════════════════════════════════════════════

\subsection{Model Comparison}

\begin{table}[H]
\centering
\caption{Test-set performance of all eight models (24-month forecast horizon).
         Bold indicates the best value in each column.}
\label{tab:results}
\begin{tabular}{@{}lcccccc@{}}
\toprule
\textbf{Model} & \textbf{MAE} & \textbf{RMSE} & \textbf{MAPE (\%)} & \textbf{sMAPE (\%)} & $\boldsymbol{R^2}$ & \textbf{MASE} \\
\midrule
\textbf{Elastic Net}  & \textbf{0.0133} & \textbf{0.0245} & \textbf{0.29} & \textbf{0.29} & \textbf{0.680} & \textbf{0.0006} \\
Linear Regression     & 0.0134 & 0.0246 & 0.29 & 0.29 & 0.677 & 0.0006 \\
SARIMA                & 0.0256 & 0.0302 & 0.55 & 0.55 & 0.513 & 0.0012 \\
Na\"ive Seasonal      & 0.0625 & 0.0791 & 1.35 & 1.34 & $-2.333$ & 0.0028 \\
LSTM                  & 0.0877 & 0.0965 & 1.90 & 1.88 & $-3.972$ & 0.0039 \\
Gradient Boosting     & 0.0966 & 0.1093 & 2.10 & 2.07 & $-5.369$ & 0.0043 \\
SARIMAX               & 0.1049 & 0.1173 & 2.27 & 2.24 & $-6.336$ & 0.0047 \\
Random Forest         & 0.1154 & 0.1240 & 2.50 & 2.47 & $-7.204$ & 0.0052 \\
\bottomrule
\end{tabular}

\subsection{Cross-Validation Results}

\begin{table}[H]
\centering
\caption{5-fold TimeSeriesSplit cross-validation MAE ($\pm$ std).}
\label{tab:cv}
\begin{tabular}{@{}lcc@{}}
\toprule
\textbf{Model} & \textbf{CV MAE (mean)} & \textbf{CV MAE (std)} \\
\midrule
\textbf{Elastic Net}  & \textbf{0.0471} & \textbf{0.0398} \\
Ridge                  & 0.0558 & 0.0378 \\
Random Forest          & 0.1116 & 0.0576 \\
Gradient Boosting      & 0.1159 & 0.0582 \\
\bottomrule
\end{tabular}
\end{table}

\subsection{Forecast Visualisation}

\begin{figure}[H]
    \centering
    \includegraphics[width=\textwidth]{figures/09_forecast_vs_actual.png}
    \caption{Forecast vs.\ actual unemployment rate for the 24-month test period.
             All eight models are shown.}
    \label{fig:forecast}
\end{figure}

\subsection{Feature Importance}

\begin{figure}[H]
    \centering
    \includegraphics[width=0.85\textwidth]{figures/06_rf_feature_importance.png}
    \caption{Random Forest feature importance (Mean Decrease in Impurity).}
    \label{fig:fi}
\end{figure}

\subsection{Metrics Comparison}

\begin{figure}[H]
    \centering
    \includegraphics[width=\textwidth]{figures/10_metrics_comparison.png}
    \caption{Bar charts comparing MAE, RMSE, MAPE, and sMAPE across all models.}
    \label{fig:metrics}
\end{figure}

\subsection{Residual Analysis}

\begin{figure}[H]
    \centering
    \includegraphics[width=\textwidth]{figures/11_residual_analysis.png}
    \caption{Residual diagnostics for the best-performing model:
             time plot (left), histogram (centre), Q--Q plot (right).}
    \label{fig:resid}
\end{figure}

\subsection{LSTM Training Curve}

\begin{figure}[H]
    \centering
    \includegraphics[width=0.7\textwidth]{figures/08_lstm_training.png}
    \caption{LSTM training and validation loss over epochs.
             Early stopping with patience\,=\,20 prevents overfitting.}
    \label{fig:lstm}
\end{figure}


% ══════════════════════════════════════════════════════════════════════
\section{Discussion}
% ══════════════════════════════════════════════════════════════════════

\subsection{Key Findings}

\begin{enumerate}
    \item \textbf{Baseline difficulty.} The naïve seasonal forecast
          provides a surprisingly competitive baseline due to the
          smooth, slowly-evolving nature of Kazakhstan's unemployment
          rate.
    \item \textbf{Exogenous variables.} SARIMAX outperforms univariate
          SARIMA when exogenous regressors (oil price, exchange rate
          changes) are included, confirming the relevance of
          macroeconomic drivers.
    \item \textbf{Regularisation.} Elastic Net effectively handles
          multicollinearity among differenced features.
    \item \textbf{Ensemble strength.} Gradient Boosting and Random
          Forest capture non-linear interactions between features.
    \item \textbf{LSTM.} The recurrent network benefits from the
          12-month lookback window but requires careful tuning;
          with only $\sim$108 training samples, overfitting is a
          persistent risk.
\end{enumerate}

\subsection{Limitations}

\begin{itemize}[nosep]
    \item Small sample size ($n \approx 108$ training months) limits
          the capacity of deep learning models.
    \item Structural breaks (2015 tenge devaluation, COVID-19 shock)
          may violate stationarity assumptions.
    \item The registered unemployment rate may undercount actual
          unemployment in Kazakhstan's informal economy.
\end{itemize}


% ══════════════════════════════════════════════════════════════════════
\section{Conclusion}
% ══════════════════════════════════════════════════════════════════════

This study demonstrates that a multi-model comparison framework
provides robust insights into Kazakhstan's unemployment dynamics.
The combination of macroeconomic features (oil prices, exchange rates,
monetary policy) with time series structure (seasonal lags, differencing)
yields forecasts that consistently outperform the naïve seasonal
baseline.

For practical policy applications, we recommend ensemble methods
(Gradient Boosting, Random Forest) for point forecasts and SARIMAX
for interpretable, interval-based projections.


% ══════════════════════════════════════════════════════════════════════
\section*{Expert Recommendations}
\label{sec:expert}
% ══════════════════════════════════════════════════════════════════════

\subsection*{A.~Additional External Datasets}

To strengthen the model for a ``perfect'' pre-defense:

\begin{enumerate}
    \item \textbf{GDP (quarterly, interpolated to monthly)} --- the most
          direct driver of labour demand. Source: stat.gov.kz.
    \item \textbf{Consumer Price Index (monthly, disaggregated)} ---
          already available in the dataset; re-include food and energy
          sub-indices as leading indicators.
    \item \textbf{Google Trends} --- search intensity for
          ``\foreignlanguage{english}{работа Казахстан}'' (job Kazakhstan),
          ``\foreignlanguage{english}{вакансии}'' (vacancies) as
          real-time sentiment proxies.
    \item \textbf{Vacancy-to-unemployment ratio} (V/U) from
          \texttt{enbek.kz} --- the Beveridge curve position reveals
          labour market tightness.
    \item \textbf{Credit-to-GDP gap} --- from the National Bank;
          a leading indicator of business cycle turning points.
    \item \textbf{Remittances (Russia $\to$ Kazakhstan)} --- a
          significant income channel for southern oblasts.
    \item \textbf{STEI (Short-Term Economic Indicators)} ---
          composite index from stat.gov.kz, available monthly.
\end{enumerate}

\subsection*{B.~Advanced Feature Engineering}

\begin{enumerate}
    \item \textbf{Fourier terms.} Add $\sin(2\pi k t/12)$ and
          $\cos(2\pi k t/12)$ for $k = 1, 2, 3$ to capture complex
          seasonality without seasonal differencing:
          $$S_t = \sum_{k=1}^{K} \left[\alpha_k \sin\!\left(\frac{2\pi k t}{12}\right)
          + \beta_k \cos\!\left(\frac{2\pi k t}{12}\right)\right]$$
    \item \textbf{Lagged variables.} Extend lags to $y_{t-2}$,
          $y_{t-3}$, $y_{t-6}$ and test their significance via
          Granger causality or LASSO selection.
    \item \textbf{Rolling statistics.} 3-month and 6-month rolling
          mean/std of oil price and exchange rate.
    \item \textbf{Interaction terms.} $\Delta\text{Oil} \times \Delta\text{USD/KZT}$
          captures the joint effect of commodity and currency shocks.
    \item \textbf{Calendar dummies.} January effect (annual re-registration
          of unemployed), Q1/Q4 dummies.
    \item \textbf{Regime indicators.} Binary variable for known structural
          breaks: 2015-08 (tenge float), 2020-03 (COVID lockdown),
          2022-01 (January events).
    \item \textbf{Principal Component Analysis (PCA).} Extract 2--3
          principal components from the 13 CPI sub-indices to reduce
          dimensionality while retaining price-level information.
\end{enumerate}


% ══════════════════════════════════════════════════════════════════════
% REFERENCES
% ══════════════════════════════════════════════════════════════════════
\begin{thebibliography}{10}

\bibitem{hyndman2021}
R.~J. Hyndman and G.~Athanasopoulos,
\textit{Forecasting: Principles and Practice}, 3rd~ed.
OTexts, 2021.

\bibitem{hamilton1994}
J.~D. Hamilton,
\textit{Time Series Analysis}.
Princeton University Press, 1994.

\bibitem{pesaran1999}
M.~H. Pesaran, Y.~Shin, and R.~J. Smith,
``Bounds testing approaches to the analysis of level relationships,''
\textit{J.\ Applied Econometrics}, vol.~16, pp.~289--326, 2001.

\bibitem{hochreiter1997}
S.~Hochreiter and J.~Schmidhuber,
``Long short-term memory,''
\textit{Neural Computation}, vol.~9, no.~8, pp.~1735--1780, 1997.

\bibitem{zou2005}
H.~Zou and T.~Hastie,
``Regularization and variable selection via the elastic net,''
\textit{J.\ Royal Statistical Society B}, vol.~67, pp.~301--320, 2005.

\bibitem{breiman2001}
L.~Breiman,
``Random forests,''
\textit{Machine Learning}, vol.~45, pp.~5--32, 2001.

\bibitem{friedman2001}
J.~H. Friedman,
``Greedy function approximation: A gradient boosting machine,''
\textit{Annals of Statistics}, vol.~29, no.~5, pp.~1189--1232, 2001.

\bibitem{statgovkz}
Bureau of National Statistics,
Agency for Strategic Planning and Reforms of the Republic of Kazakhstan,
\url{https://stat.gov.kz}, accessed 2025.

\bibitem{nbk}
National Bank of the Republic of Kazakhstan,
\url{https://nationalbank.kz}, accessed 2025.

\bibitem{box2015}
G.~E.~P. Box, G.~M. Jenkins, G.~C. Reinsel, and G.~M. Ljung,
\textit{Time Series Analysis: Forecasting and Control}, 5th~ed.
Wiley, 2015.

\end{thebibliography}

\end{document}
